\documentclass[11pt,a4paper]{report} 

% Für doppelseitigen Ausdruck (nur bei > 60 Seiten sinnvoll)
% \usepackage{ifthen}
% \setboolean{@twoside}{true}
% \setboolean{@openright}{true} 

% Deutsch
\usepackage[german]{babel} % deutsch und deutsche Rechtschreibung
\usepackage[utf8]{inputenc} % Unicode Text 
\usepackage[T1]{fontenc} % Umlaute und deutsches Trennen
\usepackage{textcomp} % Euro
\usepackage[hyphens]{url}
% statt immer Ab\-schluss\-ar\-beit zu schreiben
% einfach hier sammeln mit -. 
\hyphenation{Ab-schluss-ar-beit}
% Vorsicht bei Umlauten und Bindestrichen
\hyphenation{Ver-st\"ar-ker-aus-gang}
 % eigene Hyphenations, die für das Dokument gelten
\usepackage{amssymb} % Symbole
\usepackage{emptypage} % Wirklich leer bei leeren Seiten
\usepackage{xcolor}

%% Fonts, ein kompletter Satz an Optionen

% Times New Roman, gewohnter Font, ok tt und serifenlos
%\usepackage{mathptmx} 
%\usepackage[scaled=.95]{helvet}
%\usepackage{courier}

% Palatino mit guten Fonts für tt und serifenlos
\usepackage{mathpazo} % Palatino, mal was anderes
\usepackage[scaled=.95]{helvet}
\usepackage{courier}

% New Century Schoolbook sieht auch nett aus (macht auch tt und serifenlos)
%\usepackage{newcent}

% Oder default serifenlos mit Helvetica 
% ich kann es nicht mehr sehen ...
%\renewcommand{\familydefault}{\sfdefault}

\usepackage{microtype}

% Bilder und Listings
\usepackage{graphicx} % wir wollen Bilder einfügen
\usepackage{subfig} % Teilbilder
\usepackage{wrapfig} % vielleicht doch besser vermeiden
\usepackage{listings} % schöne Quellcode-Listings
% ein paar Einstellungen für akzeptable Listings
\lstset{basicstyle=\ttfamily, columns=[l]flexible, mathescape=true, showstringspaces=false, numbers=left, numberstyle=\tiny}
\lstset{language=python} % und nur schöne Programmiersprachen ;-)
% und eine eigene Umgebung für Listings
\usepackage{float}
\newfloat{listing}{htbp}{scl}[chapter]
\floatname{listing}{Listing}

% Seitenlayout
\usepackage[paper=a4paper,width=14cm,left=35mm,height=22cm]{geometry}
\usepackage{setspace}
\linespread{1.15}
\setlength{\parskip}{0.5em}
\setlength{\parindent}{0em} % im Deutschen Einrückung nicht üblich, leider

% Seitenmarkierungen 
\newcommand{\phv}{\fontfamily{phv}\fontseries{m}\fontsize{9}{11}\selectfont}
\usepackage{fancyhdr} % Schickere Header und Footer
\pagestyle{fancy}
\renewcommand{\chaptermark}[1]{\markboth{#1}{}}
%\fancyhead[L]{\phv \leftmark}
\fancyhead[RE,LO]{\phv \nouppercase{\leftmark}}
\fancyhead[LE,RO]{\phv \thepage}
% Unten besser auf alles Verzichten
%\fancyfoot[L]{\textsf{\small \kurztitel}}
\fancyfoot[C]{\ } % keine Seitenzahl unten
%\fancyfoot[R]{\textsf{\small Medieninformatik}}

% Theorem-Umgebungen
\newtheorem{definition}{Definition}[chapter]
\newtheorem{satz}{Satz}[chapter]
\newtheorem{lemma}[satz]{Lemma} % gleicher Zähler wie Satz
\newtheorem{theorem}{Theorem}[chapter]
\newenvironment{beweis}[1][Beweis]{\begin{trivlist}
\item[\hskip \labelsep {\textit{#1 }}]}{\end{trivlist}}
\newcommand{\qed}{\hfill \ensuremath{\square}}

% Quellen teilen
\usepackage{bibtopic} 

% Spezialpakete
\usepackage{epigraph}
\setlength{\epigraphrule}{0pt} % kein Trennstrich

% damit wir nicht so viel tippen müssen, nur für Demo 
\usepackage{blindtext} 

% Zum Zeigen von Fehlern
\usepackage{soul}
\newcommand*\falsch{\st}

 % alle Pakete und Einstellungen

% Hier anpassen 
\newcommand{\welchethesis}{Bachelor}
\newcommand{\welcherstudiengang}{Angewandte Informatik}
\newcommand{\thesisofwas}{of Science}
\newcommand{\titel}{LLM-gestützte Rekonstruktion von Softwareanforderungen
aus der Entwicklungshistorie eines Open-Source-Projekts~-- Konzeption,
Umsetzung und Evaluation am Beispiel des Quantencomputing-Frameworks
Eclipse Qrisp}
\newcommand{\kurztitel}{LLM-gestützte Anforderungsrekonstruktion}
\newcommand{\autor}{Omar Assad Salem}
\newcommand{\datum}{\textbf{[Noch zu ergänzen: Abgabedatum]}}
\newcommand{\ort}{Wiesbaden}
\newcommand{\referent}{Prof.\ Dr.\ Nikolay Tcholtchev}
\newcommand{\korreferent}{\textbf{[Noch zu ergänzen: Korreferent]}}

\begin{document}
\begin{titlepage}
  \begin{center}
    % Kopf der Seite
    \begin{tabular}{cl}
    \parbox[b]{10.2cm}{
      % \textsf würde das Aussehen der ersten Seite ruinieren, 
      % wer will, soll das selbst außen rum machen...
      \textbf{Hochschule RheinMain}\\
      Fachbereich Design Informatik Medien\\
      Studiengang \welcherstudiengang\ }
& 
\parbox[t]{6cm}{
  \includegraphics[width=0.27\textwidth]{figures/hsrmlogo.pdf}
  }
    \end{tabular}
    \vfill    
    {\LARGE \welchethesis-Arbeit} \\[0.5cm]
    {\large zur Erlangung des akademischen Grades} \\[5mm]
    {\large \welchethesis\ \thesisofwas} \\[5mm]
    \rule{\textwidth}{1pt}\\[0.5cm]
    {\begin{spacing}{1.15} \huge \bfseries \titel \\ \end{spacing}}
    \rule{\textwidth}{1pt}    
    \vfill    
    \begin{tabular}{ll} % Mitte der Seite
      Vorgelegt von & \autor \\
      am & \datum \\
      Referent & \referent \\
      Korreferent & \korreferent
    \end{tabular}    
    \vfill
  \end{center}
\end{titlepage}
\cleardoublepage

\thispagestyle{empty}
\section*{Eigenständigkeitserklärung\newline (gem. ABPO, Ziff. 4.1.5.4 bzw. RPO, §19.6)}

Ich bin verantwortlich für die Qualität des Inhalts dieser Arbeit, den ich mit geeigneten wissenschaftlichen Quellen belegt bzw. gestützt habe. Die aus fremden Quellen direkt oder indirekt übernommenen Texte, Gedankengänge, Konzepte, Grafiken, technischen Inhalte und Ähnliches in meinen Ausführungen habe ich eindeutig gekennzeichnet und mit vollständigen Verweisen auf die jeweilige Quelle versehen. Alle weiteren Inhalte der Arbeit (Textteile, Abbildungen, Tabellen etc.) ohne entsprechende Verweise stammen im urheberrechtlichen Sinn von mir.

Ich versichere außerdem, dass ich die \welchethesis-Arbeit selbständig verfasst und keine anderen als die angegebenen Quellen und Hilfsmittel benutzt sowie keine unerlaubte fremde Hilfe in Anspruch genommen habe.

%%% Anfang rote Hinweisbox zur KI-Verwendung
%%% vor Abgabe durch vereinbarte Formulierung ersetzen, Vorschläge s.u.
\begin{center}
  {\color{red}\fbox{\parbox{0.8\linewidth}{\it 
  Eine zu verwendende Formulierung zur Verwendung generativer KIs ist mit dem Referenten abzustimmen. Mögliche Formulierungen sind in der \LaTeX-Vorlage auskommentiert enthalten. Auch hiervon abweichende Regelungen sind in Absprache mit dem Referenten möglich. Bitte hier die vereinbarte Formulierung eintragen und diese rote Hinweisbox entfernen.
  }}}
\end{center}
%%% ENDE rote Hinweisbox zur KI-Verwendung

%%% Bei Verbot KI-generierter Inhalte %%%
% Ich versichere, dass ich die hier vorliegende Arbeit vollständig eigenständig erarbeitet,
% also keine KI-generierten Inhalte verwendet habe.

%%% Bei Erlaubnis KI-generierter Inhalte %%%
% Ich versichere, dass ich KI-Tools lediglich als Hilfsmittel verwendet habe und in der 
% vorliegenden Arbeit mein gestalterischer Einfluss überwiegt. Ich verantworte die Übernahme 
% jeglicher von mir verwendeter KI-generierter Inhalte vollumfänglich selbst. 

%%% (optional, falls Kennzeichnungspflicht) %%%
% Unter ``Übersicht verwendeter Hilfsmittel'' habe ich im Anhang der Arbeit sämtliche 
% KI-generierten Inhalte einzeln aufgeführt und die jeweils verwendete KI-Anwendung benannt. 
  
%%% (optional, falls Prompts abzugeben) %%%
% Die jeweiligen Benutzungseingaben (z.B. Prompts) sind der Arbeit beigefügt.


Die vorliegende Arbeit wurde bisher weder im In- noch im Ausland in gleicher oder ähnlicher Form einer anderen Prüfungsbehörde vorgelegt. Mir ist bekannt, dass ein Verstoß gegen die genannten Punkte als Täuschungsversuch gelten und prüfungsrechtliche Konsequenzen haben kann. Insbesondere kann es dazu führen, dass die Leistung nicht bestanden ist und dass bei mehrfachem oder schwerwiegendem Täuschungsversuch eine Exmatrikulation droht.

\vfill
\noindent\begin{tabular}{p{0.37\textwidth}p{0.56\textwidth}}
\ort, \datum  & \rule{0.56\textwidth}{0.5pt}\\
              & \makebox[1cm]{\ } \autor
\end{tabular}

\vfill\mbox{}
\cleardoublepage

\section*{Erklärung zur Verwendung der \welchethesis-Arbeit}

Hiermit erkläre ich mein Einverständnis mit den im folgenden 
aufgeführten Verbreitungsformen dieser Abschlussarbeit:
\bigskip

\noindent\begin{tabular}{|p{0.82\textwidth}|c|c|}
  \hline
  \textbf{Verbreitungsform} & \makebox[0.035\textwidth]{\textbf{Ja}} 
                            & \makebox[0.05\textwidth]{\textbf{Nein}} \\\hline
  Veröffentlichung des Titels der Arbeit im Internet  
                                           & $\times$ & \\\hline
  Veröffentlichung der Arbeit im Internet             
                                           &  & $\times$ \\\hline
\end{tabular}

\vspace{2cm}
\noindent\begin{tabular}{p{0.37\textwidth}p{0.56\textwidth}}
\ort, \datum  & \rule{0.56\textwidth}{0.5pt}\\
              & \makebox[1cm]{\ } \autor
\end{tabular}
\cleardoublepage



 % Titelseite, Erklärungen, etc.

\begin{abstract} 
\textbf{[Noch zu ergänzen: Abstract. Wird zuletzt geschrieben, wenn die
Ergebnisse vorliegen. Umfang etwa 150--250 Wörter: Problemstellung,
Vorgehen, zentrale Ergebnisse, Einordnung.]}
\end{abstract}

\tableofcontents

\chapter{Einleitung}
\label{cha:einleitung}

\textbf{[Noch zu ergänzen: Motivation, Problemstellung, Zielsetzung,
Forschungsfragen, Aufbau der Arbeit. Wird nach den Grundlagen
geschrieben.]}


\chapter{Grundlagen}
\label{cha:grundlagen}

\section{Qualitätskriterien für Anforderungen}
\label{sec:anforderungsqualitaet}

Die Frage, wann eine Anforderung als brauchbar gilt, ist für die
vorliegende Arbeit zentral: Ohne einen festgelegten Qualitätsmaßstab
ließe sich weder das Anforderungsformat begründet auswählen noch die
Ausgabe der Extraktionspipeline bewerten. Als Bezugsrahmen dient
ISO/IEC/IEEE~29148 in der Ausgabe von 2018, die die Anforderungsprozesse
im Lebenszyklus von Systemen und Software beschreibt~\cite{iso29148}.

\subsection{Merkmale einzelner Anforderungen}
\label{subsec:merkmale-einzeln}

Der Standard nennt neun Merkmale, die jede einzelne Anforderung
aufweisen muss. Eine Anforderung soll notwendig sein, also eine
Fähigkeit beschreiben, die ohne sie nicht abgedeckt wäre. Sie soll in
ihrer Detailtiefe zur betrachteten Abstraktionsebene passen und
Entwurfsentscheidungen nicht vorwegnehmen. Sie soll eindeutig sein und
nur eine Lesart zulassen, aus sich heraus verständlich sein und genau
einen Sachverhalt beschreiben. Weiter soll sie unter den gegebenen
Randbedingungen umsetzbar sein, ihre Erfüllung nachweisbar machen, den
zugrunde liegenden Bedarf zutreffend wiedergeben und einer festgelegten
Vorlage folgen~\cite[Abschn.~5.2.5]{iso29148}.

Für die automatisierte Rekonstruktion aus Entwicklungsartefakten sind
diese Merkmale unterschiedlich gut zugänglich. Eindeutigkeit,
Einzelbezug, Nachweisbarkeit und Vorlagenkonformität lassen sich am
Anforderungstext selbst und damit grundsätzlich auch maschinell prüfen.
Notwendigkeit, Umsetzbarkeit und Korrektheit setzen dagegen Wissen über
den tatsächlichen Bedarf und die Systemumgebung voraus, das in einer
Pull-Request-Historie nur teilweise dokumentiert ist. Diese
Unterscheidung wird in Kapitel~\ref{cha:evaluation} bei der Auswahl der
Bewertungskriterien wieder aufgegriffen.

\subsection{Merkmale einer Anforderungsmenge}
\label{subsec:merkmale-menge}

Über die einzelne Anforderung hinaus formuliert der Standard fünf
Merkmale, die sich auf die Gesamtheit der Anforderungen beziehen. Die
Menge soll vollständig sein und keine ungeklärten Platzhalter mehr
enthalten, in sich widerspruchsfrei sein und eine einheitliche
Begrifflichkeit verwenden, insgesamt umsetzbar und verständlich sein
sowie eine Validierung gegen die zugrunde liegenden Bedarfe
zulassen~\cite[Abschn.~5.2.6]{iso29148}.

Insbesondere die Forderung nach Widerspruchsfreiheit und einheitlicher
Begrifflichkeit ist für die hier entwickelte Pipeline relevant. Da die
Extraktion je Pull Request erfolgt, ist nicht ohne Weiteres
sichergestellt, dass gleichartige Sachverhalte über die gesamte
Historie hinweg gleich benannt werden. Die Behandlung von Dubletten und
Begriffsvarianten wird daher in Kapitel~\ref{cha:implementierung}
gesondert betrachtet.

\subsection{Sprachliche Kriterien}
\label{subsec:sprachkriterien}

Der Standard konkretisiert darüber hinaus, wie Anforderungen sprachlich
zu fassen sind, damit die genannten Merkmale erreicht werden.
Leitprinzip ist, dass eine Anforderung beschreibt, was benötigt wird,
und nicht, wie es umzusetzen ist. Ergänzend benennt der Standard
Klassen von Formulierungen, die zu vermeiden sind, weil sie
Mehrdeutigkeit erzeugen oder eine Verifikation verhindern: Superlative,
subjektive Bewertungen, unklare Pronomen, unbestimmte Adverbien und
Adjektive, mehrdeutige logische Verknüpfungen, offene und nicht
prüfbare Wendungen, Vergleiche ohne Bezugsgröße, einschränkende
Vorbehalte, pauschale Allaussagen sowie Verweise ohne Angabe von
Version und Datum~\cite[Abschn.~5.2.7]{iso29148}.

Diese Aufzählung ist für die vorliegende Arbeit von besonderem Wert, da
sie sich in maschinell prüfbare Regeln überführen lässt. Sie bildet
damit die Grundlage sowohl für die automatische Validierung der
Modellausgaben als auch für ein quantitatives Vergleichsmaß im
Modellvergleich in Kapitel~\ref{cha:evaluation}.

\section{Anforderungsformate}
\label{sec:anforderungsformate}

\textbf{[Noch zu ergänzen: RFC 2119, User Stories, EARS, die
Spezifikationsvorlagen aus ISO/IEC/IEEE 29148.]}

\section{Large Language Models und Prompt Engineering}
\label{sec:llm-grundlagen}

\textbf{[Noch zu ergänzen]}

\section{Common Criteria und BSI-orientierte Prüfungen im Überblick}
\label{sec:cc-bsi}

\textbf{[Noch zu ergänzen]}


\chapter{Verwandte Arbeiten}
\label{cha:verwandte-arbeiten}

\textbf{[Noch zu ergänzen: LLMs im Requirements Engineering,
retrospektives Requirements Engineering, Mining von
Entwicklungsartefakten.]}


\chapter{Konzeption der Pipeline}
\label{cha:konzeption}

\textbf{[Noch zu ergänzen: Formatauswahl, Architektur, Datenmodell,
Prompt-Design, Klassifikationsraster, Modellwahl,
Evaluationskonzept.]}


\chapter{Implementierung}
\label{cha:implementierung}

\textbf{[Noch zu ergänzen: GitHub-API, Aufbereitung, Extraktion,
Klassifikation, Gap-Analyse, Absicherung.]}


\chapter{Evaluation}
\label{cha:evaluation}

\textbf{[Noch zu ergänzen: Design, Durchführung, Ergebnisse,
Auswertung.]}


\chapter{Diskussion}
\label{cha:diskussion}

\textbf{[Noch zu ergänzen: Beantwortung der Forschungsfragen,
Interpretation, Übertragbarkeit, Validitätsbedrohungen,
Limitationen.]}


\chapter{Fazit und Ausblick}
\label{cha:fazit}

\textbf{[Noch zu ergänzen]}


\bibliographystyle{plain} % Literaturverzeichnis
\begin{btSect}{thesis} % mit bibtopic Quellen trennen
\section*{Literaturverzeichnis}
\btPrintCited
\end{btSect}
\begin{btSect}{online}
\section*{Online-Quellen}
\btPrintCited
\end{btSect}

\end{document}
